\chapter{Module categories, exact MPS, and gapped alternatives}
\label{ch:gapped}

An unusual noninvertible symmetry can restrict the possible vacuum
structure without forcing a gapless theory. To see this concretely, we
construct finite-bond states from a Haagerup algebra object and explain
why a sufficiently long-range parent Hamiltonian is gapped and symmetric
on the same physical-rho path space as the target chain. We then derive
the corresponding finite-dimensional vacuum form factors. These are
valuable alternatives against which to compare the chain data, but they
are not assertions that the target Hamiltonian realizes that phase.
The needed categorical ideas are algebra objects, their modules, and
the compatibility of an action with reassociation.
The matrix-product notation comes from Chapter~\ref{ch:tensor_networks};
the relation between algebra modules and physical defects was introduced
in Chapter~\ref{ch:defects}.

\section{Why several vacua can be required}
\index{module category!gapped vacua}

For an ordinary finite symmetry group, a symmetry-broken gapped phase
has a finite set of pure vacua permuted by the group. A noninvertible
line can send one vacuum to a sum of vacua. Its multiplicities form
nonnegative integer matrices satisfying the fusion rules. These matrices
are often called a nonnegative integer matrix representation, or NIMrep.

The matrices are not all the data. There must also be consistent
associativity maps describing how line junctions act on interfaces
between vacua. Those maps make the vacua into a module category over
the symmetry category. A representation of the abstract fusion ring on
a Hilbert space need not lift to such a positive semisimple vacuum
module. This distinction prevents a small matrix trick from proving the
existence of a physical phase.

An invariant finite-volume ground state can be a cat of several pure
vacua. Its one-point functions can vanish even when each component is
ordered. In an ordinary gapped tunneling picture, several finite-volume
levels collapse toward a common vacuum energy while a separate bulk gap
remains positive. Their splittings may be exponentially small at large
circumference. A broad crossover can make those splittings look like
$1/L$ before the asymptotic regime is reached.

\begin{example}[A gapped ordered chain on compatible and frustrated rings]
\label{ex:gapped-clock-frustration}
\index{clock model!period-three order}\index{domain wall}
Consider an ordinary three-state clock at each site, with orthonormal
basis $|s_j\rangle$, $s_j\in\Z/3\Z$. Define a diagonal Hamiltonian
by its energy on each product configuration:
\[
 H_{\rm clock}|s_1,\ldots,s_L\rangle
 =J\sum_{j=1}^L\bigl(1-\delta_{s_{j+1},\,s_j+1}\bigr)
       |s_1,\ldots,s_L\rangle,
 \qquad s_{L+1}=s_1,\quad J>0.
\]
Take $L\ge3$. All color arithmetic is modulo three. The local
terms are commuting projections times $J$: a satisfied bond
advances the color by one, while an unsatisfied bond costs $J$.
The three zero-defect configurations on the infinite chain are
$s_j=a+j$, with $a\in\Z/3\Z$. These are period-three pure bulk
vacua, cycled by one-site translation or a global color shift.

The finite spectrum follows from the defect residues
\[
 r_j=s_{j+1}-s_j-1\in\Z/3\Z,
 \qquad \sum_{j=1}^L r_j=-L\pmod3.
\]
Their nonzero entries count the energy in units of $J$.
Conversely, any residue sequence satisfying this sum, together
with one initial color, reconstructs a unique periodic configuration.
When $3\mid L$, all $r_j=0$ is allowed and gives exactly three
ground states of energy zero. One nonzero residue cannot have
sum zero. Two residues $+1,-1$ can, so the first excited energy
is exactly $2J$. With energies ordered from $E_0$ and counted
with multiplicity,
\[
 E_0=E_1=E_2=0,\qquad E_3-E_0=2J
 \quad(3\mid L).
\]
There is a vacuum cluster and a size-independent gap above it.

If $3\nmid L$, at least one defect is forced. One residue
$r_j=-L$ satisfies the constraint; there are $L$ choices of
its bond and three choices of initial color. Hence the ground
energy is $J$ with degeneracy $3L$. In particular $E_3-E_0=0$
along these sizes. This does not diagnose a critical bulk:
the incompatible ring contains a domain wall whose position
costs no energy in this exactly diagonal model.
The staggered variable $t_j=s_j-j$ makes the mechanism transparent.
In the bulk the Hamiltonian is the ferromagnetic clock projector
onto unequal neighboring $t$ values; the ring closes with
$t_{L+1}=t_1-L$, a nontrivial color twist precisely when
$3\nmid L$.

The distinction between a pure vacuum and a cat is equally explicit.
Let $Z_j|s_j\rangle=\omega^{s_j}|s_j\rangle$, with
$\omega=e^{2\pi i/3}$. In each product vacuum connected
correlations between distinct sites vanish. On a compatible ring,
the equal superposition of the three product vacua instead has
\[
 \langle Z_j\rangle=0,\qquad
 \langle Z_i^\dagger Z_j\rangle=\omega^{j-i}.
\]
Its nondecaying order correlation coexists with the gap $2J$.
This elementary clock model illustrates the finite-volume issues;
it makes no assertion about the phase of the Haagerup Hamiltonian.
\end{example}

The physical $\mathcal H_3$ category has indecomposable vacuum module possibilities
of ranks two, four, and six. Translation can also permute equivalent
components. Thus the word ``symmetric'' need not mean a unique pure
vacuum. A finite collection of vacua can carry the full symmetry
consistently.

\section{An algebra object is a multiplication with \texorpdfstring{$F$}{F}-moves}
\index{Q-system!Haagerup module algebra}

An algebra object consists of an object $Q$, a unit
$\eta:1\to Q$, and a multiplication $m:Q\otimes Q\to Q$ satisfying
unit and associativity equations. Associativity compares
$m(m\otimes1)$ with $m(1\otimes m)$ after applying the category's
associator. In a fusion basis it is a system of polynomial equations
whose coefficients are the $F$-symbols.

For a unitary separable Frobenius algebra, the adjoint $m^\dagger$
is compatible with multiplication, and a normalization can make
\[
 mm^\dagger=\qdim(Q)\,1_Q.
\]
The Frobenius identities are not decorative. They are what lets
multiplication and comultiplication pass through one another in a
positive tensor construction. A connected algebra has
$\dim\Hom(1,Q)=1$.

\begin{example}[The associator can forbid an apparent algebra]
One might try $Q'=1+\rho$, since its summands look simple and its
fusion products contain the needed output channels. Normalize the
unit and write
$u=m(\rho,\rho;1)$, $v=m(\rho,\rho;\rho)$. In the physical $\mathcal H_3$
associator, the two forbidden noninvertible output channels in the
three-rho associativity equation require
\[
 \frac u{\sqrt d}+y_1v^2=0,\qquad
 \frac u{\sqrt d}+y_2v^2=0,
\]
where $y_1\ne y_2$. Hence $u=v=0$. The remaining multiplication is
a square-zero extension; its pairing through the tensor-unit channel
is degenerate, so it is not a connected separable unitary algebra.
Fusion-ring closure alone therefore fails to construct $1+\rho$ as
a Q-system for this associator.
\end{example}

This local calculation also distinguishes the physical rank-four
$\mathcal H_3$ module from a superficially similar $\mathcal H_2$ module. A category name
or a shared fusion graph is not a substitute for checking its
associativity equations.

\begin{constructioninput}[The physical Haagerup algebra]
In the specified real $\mathcal H_3$ gauge, the repository constructs
\[
 Q=1+\rho+\alpha\rho
\]
with unit, associativity, Frobenius, and specialness identities
checked exactly in its algebraic coefficient field. It also constructs
the four simple right modules and their action coefficients. The
following MPS and parent arguments take these finite algebraic
identities as input; they do not infer them from numerical
optimization.
\end{constructioninput}

The finite input can be made quite explicit. Put $r=\rho$ and
$s=\alpha\rho$, and define real core constants
\[
 x=\frac{2-\sqrt{13}}3,\qquad
 y_{1,2}=\frac{5-\sqrt{13}\mp\sqrt{6(1+\sqrt{13})}}{12},
 \qquad z=\frac{1+\sqrt{13}}6.
\]
Here the minus sign defines $y_1$. Set
\[
 D=\sqrt{z^2-y_1^2},\quad u=\sqrt d,\quad
 v=-y_1/D,\quad\zeta=z/D,\quad
 w=\sqrt{d-1-v^2},\quad t=-\frac{y_2}{z}w,
\]
with positive square roots. The unit coefficients are one, and the
remaining nonzero multiplication coefficients are
\begin{align*}
 m(r,r;1)&=m(s,s;1)=u,\\
 m(r,r;r)&=v,\\
 m(r,r;s)&=m(r,s;r)=m(s,r;r)=w,\\
 m(s,s;r)&=m(r,s;s)=m(s,r;s)=\zeta,\\
 m(s,s;s)&=t.
\end{align*}
These are coefficients on normalized trivalent fusion vertices, not
the multiplication table of an ordinary three-dimensional vector
space with a trivial associator. Every unlisted coefficient is zero.

\begin{example}[Two specialness blocks]
On the tensor-unit output, $mm^\dagger$ is the sum of the squares
of $m(1,1;1)$, $m(r,r;1)$, and $m(s,s;1)$, namely
$1+2u^2=1+2d=\qdim(Q)$. On the $r$ output, it is
\[
 2+v^2+2w^2+\zeta^2.
\]
The radical definitions give $\zeta^2-v^2=1$ and
$w^2=d-1-v^2$, reducing the expression to the same $1+2d$.
These checks illustrate how the coefficients are normalized. They
do not replace the associativity equations, including forbidden
channels, or the separate Frobenius checks.
\end{example}

In the convention of \eqref{eq:chain-F-convention}, a full
associativity component has the concrete form
\[
 \sum_e m(a,b;e)m(e,c;T)F^{abc}_T[e,f]
       =m(b,c;f)m(a,f;T).
\]
All allowed $f$ and all simple totals $T$, including those outside
$Q$, enter the test. This is a small enough example to see exactly
what an algebra-object construction claims while still using a
nontrivial Haagerup associator.

A right module is an object $M$ with action
$\mu_M:M\otimes Q\to M$ satisfying the analogous unit and
associativity equations. The category acts on modules from the left:
$c\otimes M$ inherits a right action by reassociation. Decomposing
this induced module into simples gives the vacuum fusion matrices.
Here $Q$ belongs to the nonbraided line category $\mathcal H_3$.
Its right modules form the category of vacua on which the lines act;
the argument requires no independent tensor product of two vacua.

\section{The four right modules and their physical tensors}

Three modules are $\alpha^jQ$, with underlying supports
\[
 \alpha^j+\alpha^j\rho+\alpha^{j+1}\rho.
\]
The fourth, denoted $M_{\rm fix}$, has underlying object
$\rho+\alpha\rho+\alpha^2\rho$ and is fixed by the alpha action
up to module isomorphism. Its existence can be seen from the free
module $\rho\otimes Q$ after removing two known simple summands.

The multiplicity vector of $\rho Q$ in the ordered six-object list
is $(1,0,1,3,2,2)$. It decomposes as
\[
 \rho Q\cong Q\oplus\alpha^2Q\oplus M_{\rm fix}.
\]
The exact action coefficients provide orthogonal coisometries onto
the first two summands; their orthogonal complement gives the last
one. The endomorphism dimension
$\dim\End_Q(\rho Q)=\dim\Hom(\rho,\rho Q)=3$ then shows that the
nonzero complement is simple, once the first two distinct simple
summands have been identified.

Restrict the physical algebra input of a module action to the
$\rho$ summand of $Q$. In the multiplicity-free modules here, this
gives scalar edge weights
\[
 W_M(a,b)=\mu_M(a,\rho;b).
\]
The periodic state is
\begin{equation}
 \Psi_M(a_0,\ldots,a_{L-1})
 =\prod_iW_M(a_i,a_{i+1}).
 \label{eq:module-edge-state}
\end{equation}
More generally one takes an ordinary trace over multiplicity fibers.
Forbidden fusion edges vanish, so this is an exactly admissible
physical-rho state.

\subsection{Why the categorical line has the expected action}

Apply the $F$-product operator $Y_c$ from
\eqref{eq:topological-line-product} to \eqref{eq:module-edge-state}.
At each output height, summing the internal input height is precisely
the induced right action on $c\otimes M$. A heightwise module
decomposition turns this action into a direct sum of simple module
actions. The changes of basis cancel on consecutive edges and inside
the final cyclic trace. Therefore
\begin{equation}
 Y_c\Psi_M=\sum_NN_{c,M}^{N}\Psi_N.
 \label{eq:module-line-closure}
\end{equation}
No quantum-dimension prefactor appears. The closure uses ordinary
cyclic traces in the orthonormal fusion coordinates; a categorical
pivotal trace would be a different contraction.

The finite constructor checked this statement literally on every
periodic path at lengths three and four, as well as checking the
module equations from which it follows generally. Those small-ring
checks validate the implementation and normalization. The all-length
argument is the telescoping module decomposition, not an extrapolation
from two rings.

After a cyclic relabeling of the three ordered modules, the rho
multiplicity matrix is
\begin{equation}
 R=\begin{pmatrix}
 1&1&1&1\\
 1&0&1&1\\
 1&1&1&0\\
 1&1&0&1
 \end{pmatrix}.
 \label{eq:H3-rank4-R}
\end{equation}
Alpha fixes the first state and cycles the last three. Its triplet
block is $J-P$, where $P$ fixes one color and exchanges the other
two. This is the category-correct rank-four $\mathcal H_3$ action. In
\eqref{eq:module-line-closure} it acts on the canonical, generally
unnormalized finite-ring module tensors. Those tensors need not be
orthogonal at finite circumference. In the gapped vacuum description
below, the same symmetric matrix acts on an orthonormal basis of
simple vacua. These are distinct normalization statements; independently
normalizing four finite-ring vectors would insert ratios of their norms
into the matrix elements.

\section{A scalar-edge family that can be solved by hand}

For the alpha-fixed module, a real height gauge makes all diagonal
weights equal and all off-diagonal weights equal. After an overall
normalization, consider the larger one-parameter family
\[
 W_t(i,j)=\begin{cases}t,&i=j,\\1,&i\ne j,\end{cases}
 \qquad i,j\in\Z/3\Z,\quad t\in\R.
\]
Its norm transfer is
\[
 K_t=(t^2-1)I+J.
\]
The leading eigenvalue is $t^2+2$ and the two other eigenvalues are
$t^2-1$. The leading probability on colors is uniform. Even at $t=0$
the transfer is primitive: a complete graph on three vertices without
loops has both two-step and three-step returns and hence period one.
For nonzero $t$, the positive diagonal makes aperiodicity immediate.

The module point is a particular negative solution
\[
 t_0^2+t_0=1/d.
\]
The nonzero amplitudes on all three colors show how an alpha-invariant
primitive state can live entirely on noninvertible heights. There is
no requirement that every height appear with nonzero probability in
every allowed physical state.

The energy of this entire family is also a small local contraction.
Write $x,y_1,y_2,z$ for the relevant rho-projector column entries in
the physical $F$ gauge. For equal endpoint colors, the two-edge
amplitude vector is $(t^2,1,1)$ and its rho-channel component is
$xt^2+y_1+y_2$. For unequal endpoint colors it is a permutation of
$(t,t,1)$ and its component is $(y_1+y_2)t+z$. The transfer Perron
normalization gives
\begin{equation}
 e(t)=-\frac{(xt^2+y_1+y_2)^2+
                   2((y_1+y_2)t+z)^2}{(t^2+2)^2}.
 \label{eq:scalar-edge-energy}
\end{equation}
This is an exact infinite-state variational energy, not a finite-size
fit. At the exact negative module point it is only about $-.065$,
far above the optimized physical-chain energy density near $-.845$.
The construction supplies a controlled state and a symmetry example,
not a competitive ground-state approximation by itself.

Module associativity gives another exact property. Recoupling two
adjacent physical rho legs into a channel $f$ produces
\[
 m(\rho,\rho;f)\,\mu_M(a,f;b).
\]
Since $\alpha^2\rho$ is absent from $Q$, its coefficient vanishes.
Every constructed module state is therefore annihilated by every
$P_i^{(\alpha^2\rho)}$. It is a frustration-free state of the
different Hamiltonian $+\sum_iP_i^{(\alpha^2\rho)}$. It has not
thereby been shown that this short Hamiltonian has no other ground
states or has a positive thermodynamic gap.

\section{A gapped parent on the same constrained path space}
\index{parent Hamiltonian}

We now explain the stronger long-range parent construction. Its
input is the four exact tensors, and its conclusion concerns a
particular sufficiently long interaction range. It does not change
the physical object per site.

For an ordered block $Q$, in the support order $(1,\rho,\alpha\rho)$,
the weight matrix has the form
\[
 W_Q=\begin{pmatrix}0&1&0\\u&v&w\\0&w&\zeta\end{pmatrix},
 \qquad u,v,w,\zeta>0.
\]
It is invertible because $\det W_Q=-u\zeta$. Its norm graph has
links $I\leftrightarrow R\leftrightarrow S$ and loops on $R,S$,
so it is primitive. In the real height gauge, the fixed block is
$W_f=B_f[(t_0-1)I+J]$, where $B_f>0$ is its common off-diagonal
weight. It is also invertible:
\[
 \det W_f=B_f^3(t_0+2)(t_0-1)^2\ne0.
\]
Neither $t_0=-2$ nor $t_0=1$ solves $t_0^2+t_0=1/d$.

Embed each block as an ordinary three-dimensional virtual tensor
\[
 A_M^a=|a\rangle\operatorname{row}_a(W_M),
\]
with zero tensor when $a$ is outside the support. A physical word
then yields
\[
 A_M^{a_0}\cdots A_M^{a_{\ell-1}}
 =\left[\prod_{j=0}^{\ell-2}W_M(a_j,a_{j+1})\right]
 |a_0\rangle\operatorname{row}_{a_{\ell-1}}(W_M).
\]
The products of words span the virtual matrix algebra once enough
words distinguish the blocks and vary their endpoints.

\begin{proposition}[An explicit block-injectivity witness]
Five-site words span the direct sum of the four $3\times3$ virtual
matrix algebras of these module tensors.
\end{proposition}
\begin{proof}
In ordered block $j$, choose a word
$a,R_j,I_j,R_j,b$ with arbitrary supported endpoints. Every edge
has nonzero weight. The unique invertible height $I_j$ excludes
the other three modules. Varying $a,b$, and using invertibility of
$W_j$, spans all nine matrix units in that block.

For the fixed block, choose an all-noninvertible five-site word with
any prescribed endpoints that visits all three colors. Every edge
is nonzero. No ordered module supports all three colors, so it
isolates the fixed block. Again the endpoint matrices span all
nine matrix units.\end{proof}

Block injectivity is stronger than merely finding a primitive
transfer spectrum separately in each block. It also separates the
blocks by physical words. It is the appropriate input for a
finite-family MPS parent theorem.
Equivalently, after blocking, every virtual matrix can be detected
by a suitable physical word within its chosen block. This is the
linear-algebra meaning of \emph{injectivity} being used here.

For an interval of length $\ell$, define the local range
\[
 S_\ell=\Span\{|v^M_{ab}\rangle:M,a,b\},
\]
where $|v^M_{ab}\rangle$ sums all internal height words with fixed
endpoints $a,b$, weighted by the product of the $W_M$ edges. Because
$W_M$ is invertible, this is precisely the ordinary MPS range with
arbitrary virtual boundary matrices. Different endpoint pairs have
orthogonal physical support. Its orthogonal projector $\Pi_\ell$
therefore preserves both endpoint heights and maps admissible
interiors to admissible interiors.

The parent term $h_\ell=1-\Pi_\ell$ preserves the global constrained
path space. In the unconstrained six-state embedding, forbidden
interior words are penalized. Restricting to the path space is a
reducing restriction of the operator, not a change in its intended
low-energy state family.

\subsection{Why the local parent respects the categorical lines}

Closure of the four periodic states under $Y_c$ would not by itself
prove that every local parent term commutes with $Y_c$. The needed
statement is local. Fix the endpoints of an interval and all outside
MPO indices. Acting with the open $F$-product line on an interval
module vector is the induced module $c\otimes M$ on that interval.
The same module decomposition used in
\eqref{eq:module-line-closure} telescopes at every internal height,
leaving only boundary intertwiners. Hence it sends the whole local
range $S_\ell$ into itself.

Thus the subspace with this interval in $S_\ell$ and arbitrary
compatible outside data is invariant under every closed $Y_c$ and
under $Y_c^\dagger=Y_{c^*}$. It is therefore reducing. Its
orthogonal projector commutes with the closed line on the full
compatible chain, giving $[h_\ell,Y_c]=0$.

For a sufficiently long fixed interval, the standard block-injective
MPS parent theorem gives a thermodynamic gap and four periodic
ground states on sufficiently large rings. All four MPS lie in the
admissible path space. Restricting the positive parent to this
reducing subspace keeps those ground states and cannot create
eigenvalues below its gap. The same local term can
be placed at every site, preserving one-site translation. The
five-site word argument is a proof device, not a requirement to
break physical translation to a five-site cell.

\begin{theorem}[Conditional on the exact module input]
The physical one-rho-per-site path Hilbert space admits a finite-range,
one-site-translation-invariant gapped Hamiltonian commuting with the
actual $\mathcal H_3$ categorical lines and realizing the rank-four module
of vacua.
\end{theorem}

No minimal interaction range or numerical gap bound is asserted.
In particular, this theorem does not prove a gap for the short
missing-channel anchor. Its open kernels exceed the module ranges
already at small lengths. Nor does it determine the phase of
$-\sum P_\rho$. It rules out a proposed kinematic argument that
these symmetries and this object per cell force every compatible
Hamiltonian to be gapless.

\begin{example}[A local kernel does not determine the longer parent]
The distinction can be checked by exact dimensions. On three height
sites, the four-module range is in fact the entire kernel of the
single $P_{\alpha^2\rho}$ term. At equal noninvertible endpoints,
there are four middle heights and a rank-one constraint. The three
available module vectors are independent: one has an invertible
middle height unavailable to the others; after removing it, the
fixed module has a third-color middle height unavailable to the
remaining ordered module. At unequal noninvertible endpoints, the
middle space has dimension three and the two module vectors are
independent, since only the fixed module reaches the third color.

The other endpoints contribute one dimension whenever they belong
to an ordered module. Summing equal invertible, mixed, equal
noninvertible, and unequal noninvertible endpoint sectors gives
\[
 \dim S_3=3+6+6+9+12=36.
\]
Thus three-site words already separate the four virtual blocks;
the five-site witness above was a simpler sufficient proof. The
sum of the physical tensors is the invertible direct sum of their
$W_M$ matrices, so full spanning persists at every larger word
length. In particular $\dim S_4=36$.

Now let $D_\ell$ count admissible open words of $\ell$ heights,
including all endpoints. Unlike the periodic count, it is the sum
of all entries of $N_\rho^{\ell-1}$. The color-independent
two-by-two block gives $D_1=6$, $D_2=15$ and
$D_{\ell+2}=3D_{\ell+1}+D_\ell$, hence
\[
 D_3=51,\qquad D_4=168.
\]
Each of the two missing-channel projectors on four heights has
rank $51$: fusing the selected pair into $\alpha^2\rho$ reduces
the edge count by one, and its alpha permutation disappears when
all endpoints are summed. Rank subadditivity therefore gives
\[
 \dim\bigl(\ker P_1\cap\ker P_2\bigr)
 \geq168-51-51=66>36=\dim S_4.
\]
No commutativity assumption is used. The equality of one local
kernel with an MPS range has already failed to propagate at the
next length. This is why a parent-gap theorem must use its actual
intersection hypotheses, rather than only the existence of several
annihilated MPS.
\end{example}

\section{The rank-four vacuum algebra and its exact form factors}
\index{vacuum algebra}

We can derive a sharp low-energy alternative if the target chain
eventually approaches this gapped vacuum structure. Let $|f\rangle$
be the fixed pure vacuum and $|o_j\rangle$ the ordered three-cycle.
Put
\[
 b=d-2=\frac{\sqrt{13}-1}2,\qquad b^2=3-b.
\]
Choose the vacuum in the trivial categorical sector,
$Y_\alpha=1$ and $Y_\rho=d$. The four-dimensional ground space
alone does not make this choice; alpha invariance still leaves a
two-dimensional neutral space. The Perron eigenvector of
\eqref{eq:H3-rank4-R} with eigenvalue
$d$ is
\[
 |0\rangle=\frac{b|f\rangle+|o_0\rangle+|o_1\rangle+|o_2\rangle}
                  {\sqrt{b^2+3}}.
\]
The ordered total probability is $p=3/(b^2+3)$. The symmetric
finite-volume vacuum is therefore not an equal mixture of four
vacua.

There is a useful distinction between state normalization and
idempotent normalization here. Let $v_f,v_j$ be the local primitive
idempotents, which multiply as $v_i v_j=\delta_{ij}v_i$, and let
$\theta_i=\langle v_i,v_i\rangle$ in the positive Frobenius
inner product. The line action on these idempotents is
\[
 Y_\rho v_j=(J-P)v_j+b^{-1}v_f,
 \qquad Y_\rho v_f=v_f+b\sum_jv_j.
\]
Self-adjointness gives $\theta_f=b^2\theta_o$, with the three
ordered norms equal. Passing to $|i\rangle=v_i/\sqrt{\theta_i}$
produces the symmetric integer matrix
\eqref{eq:H3-rank4-R}. Thus the unequal Perron probabilities are
fixed by the positive vacuum metric; an equal-weight statistical
mixture would calculate different form factors.

In the diagonal algebra of functions on these four vacua, define
\[
 \Phi(o_j)=p^{-1/2}\omega^j,\quad \Phi(f)=0,
 \qquad E(o_j)=b/\sqrt3,\quad E(f)=-\sqrt3/b.
\]
Direct averaging in $|0\rangle$ gives
$\langle E\rangle=0$ and
$\langle E^2\rangle=\langle\Phi^\dagger\Phi\rangle=1$.
The cubic and mixed form factors are
\begin{align}
 C_{\Phi\Phi\Phi}&=p^{-1/2}
       =\sqrt{\frac{13-\sqrt{13}}6},\label{eq:rank4-charged}\\
 g=C_{\Phi\Phi^\dagger E}&=b/\sqrt3,
 \qquad C_{EEE}=-1/\sqrt3.\label{eq:rank4-neutral}
\end{align}
For instance, $\Phi^3$ is $p^{-3/2}$ on the ordered vacua, so its
mean is $p^{-1/2}$. The neutral cubic follows by adding the ordered
and fixed contributions with their unequal probabilities.

The algebra closes:
\begin{equation}
 \Phi^\dagger\Phi=1+gE,\qquad
 E^2=1+C_{EEE}E,\qquad
 |C_{\Phi\Phi\Phi}|^2=1+g^2,
 \quad C_{EEE}=\frac{g^2-1}{g}.
 \label{eq:rank4-vacuum-identities}
\end{equation}
These are ordinary multiplication identities in a finite-dimensional
vacuum algebra. They do not use conformal invariance. A microscopic
operator ratio approaches them only if its source and selected states
converge to the stated vacuum projections.

The physical $\mathcal H_3$ predictions are approximately $1.25$ for the
charged cubic and $.75$ for the mixed neutral coefficient. The
values about $1.66$ and $1.33$ arise from the different rank-four
$\mathcal H_2$ action. The finite-chain values near $1.59$ and $1.30$ differ
from the $\mathcal H_3$ targets, but do not alone exclude a late gapped
crossover: vacuum projection and a finite bulk correlation length
have not been established at these sizes.

\section{Other vacuum modules and a symmetry filter}

In the rank-two module, alpha is trivial and
$R=\left(\begin{smallmatrix}0&1\\1&3\end{smallmatrix}\right)$.
Its Perron amplitudes are proportional to $(1,d)$. The centered
unit neutral field has values $d,-d^{-1}$ and cubic
$d-d^{-1}=3$. There is no charged order field in this vacuum
algebra.

The regular rank-six module has two alpha triplets, with Perron
amplitudes $(1,1,1,d,d,d)$. It has two charged pairs rather than
one, and a neutral cubic of $3$. Its charged cubic depends on the
chosen linear combination of the two pairs. Under an additional
closed-pair condition for the rho line, its upper possible magnitude
is close to $1.5895$. Matching that one number is not an
identification: the extra charged pair and the distinct neutral
predictions must also occur.

For the rank-four module, let
$|s\rangle=(|o_0\rangle+|o_1\rangle+|o_2\rangle)/\sqrt3$.
On $(f,s)$,
\[
 Y_\rho=\begin{pmatrix}1&\sqrt3\\\sqrt3&2\end{pmatrix},
 \qquad (Y_\rho-2)|s\rangle=\sqrt3|f\rangle.
\]
This suggests extracting the missing fixed vacuum from an ordered
cat. On the whole alpha-neutral Hilbert space, however,
$B=Y_\rho-2$ is invertible, not a projection, and
$B^\dagger B=5-Y_\rho$. For a generic normalized state $s$,
\[
 \|Bs\|^2=5-\langle Y_\rho\rangle.
\]
Dividing by $\sqrt3$ is correct only for the exact module cat with
$\langle Y_\rho\rangle=2$.

Since $[H,Y_\rho]=0$, the normalized filtered energy obeys
\[
 E_B-E=-\frac{\langle HY_\rho\rangle-E\langle Y_\rho\rangle}
                 {5-\langle Y_\rho\rangle}.
\]
If $\operatorname{Var}(H_L)=O(L)$, its energy-density difference
vanishes at least as $L^{-1/2}$ by Cauchy--Schwarz. Equal limiting
energy densities after this symmetry filter are therefore a general
variance consequence, not proof of a new pure vacuum.

A coherent tensor subtraction must retain mixed overlaps. Taking
separate Perron limits of $Y_\rho s$ and $s$ before subtracting can
destroy the cancellation. At finite bond dimension, two close but
inequivalent injective tensors can become orthogonal as the ring
length tends to infinity. A failed finite-bond extraction therefore
does not refute the rank-four phase, while small charged means do
not prove primitivity.

\section{Exercises}

\begin{exercise}\label{ex:gapped-PF}
Verify the fusion relation for \eqref{eq:H3-rank4-R} together with
the alpha cycle. Find its eigenvalues on the neutral two-plane and
its action on the charged complement.

\par\smallskip\noindent{\footnotesize\normalfont
\hyperref[sol:ex:gapped-PF]{Solution on page~\pageref*{sol:ex:gapped-PF}.}}\par
\end{exercise}
\begin{exercise}\label{ex:gapped-cubics}
Derive all three coefficients in
\eqref{eq:rank4-charged}--\eqref{eq:rank4-neutral} and the identities
\eqref{eq:rank4-vacuum-identities} by averaging functions on four
points with the Perron probabilities.

\par\smallskip\noindent{\footnotesize\normalfont
\hyperref[sol:ex:gapped-cubics]{Solution on page~\pageref*{sol:ex:gapped-cubics}.}}\par
\end{exercise}
\begin{exercise}\label{ex:gapped-parent}
Explain why invariance of the four periodic MPS under a closed line
does not by itself imply symmetry of each local parent projector.
Identify the stronger open-interval property used in the proof.

\par\smallskip\noindent{\footnotesize\normalfont
\hyperref[sol:ex:gapped-parent]{Solution on page~\pageref*{sol:ex:gapped-parent}.}}\par
\end{exercise}
\begin{exercise}\label{ex:gapped-filter}
Derive the filtered-energy identity and its $O(L^{-1/2})$ density
bound when $\operatorname{Var}(H_L)=O(L)$. Give an example of a
cat of states with different energy densities for which the variance
is instead of order $L^2$.

\par\smallskip\noindent{\footnotesize\normalfont
\hyperref[sol:ex:gapped-filter]{Solution on page~\pageref*{sol:ex:gapped-filter}.}}\par
\end{exercise}
\begin{exercise}\label{ex:gapped-short}
Suppose several MPS are annihilated by every term of a short
frustration-free Hamiltonian. List the two additional conclusions
needed before one can claim this Hamiltonian has exactly those
ground states and a thermodynamic gap. Why does a gapped longer
parent not supply them automatically?

\par\smallskip\noindent{\footnotesize\normalfont
\hyperref[sol:ex:gapped-short]{Solution on page~\pageref*{sol:ex:gapped-short}.}}\par
\end{exercise}
